% ============================================================================
% Proofs appendix for sections section_germ.tex (§2) and section_geometry.tex (§3).
% Same environment
% requirements. Numbering: statements are restated with \ref to the main text.
% ============================================================================

\section{Proofs}
\label{app:proofs}

\paragraph{Verification protocol.} Every identity proved below was also
verified numerically at float64 by two independent adversarial verification
agents with separately written code (including an independent implementation
of the segment traversal for Theorem~\ref{thm:anatomy}); worst-case measured
errors are quoted where informative. Scripts and outputs accompany the
submission package.

\subsection{Preliminaries}

\begin{clarifybox}
\textbf{The vocabulary of this lemma, spelled out.} Fix a network satisfying
(LCS), so that by Proposition~\ref{prop:lcs}(ii) each activation has the form
$f(z)=a_+\max(z,0)+a_-\min(z,0)$.

\emph{Pattern, and $\sigma$.} The \emph{pattern} at $x$ is the assignment to
each hidden unit $(\ell,q)$ of the side of the origin its pre-activation falls
on, equivalently of which of the two slopes $a_\pm$ its entry of $D^{(\ell)}(x)$
takes. It is a vector of bits, one per hidden unit --- for ReLU exactly the
familiar $0/1$ mask. A \emph{candidate} pattern is written $\sigma$: any one of
the $2^{\#\mathrm{units}}$ bit assignments, whether or not some input realises
it. There are finitely many, which is what makes the union below finite.

\emph{Where the activation ``breaks''.} A \emph{breakpoint} of $f$ is a point
where its slope changes, i.e.\ where the two linear pieces meet. Under (LCS)
the only breakpoint is $z=0$ (Proposition~\ref{prop:lcs}(ii)): away from $z=0$
the function is linear and $D$ is locally constant, so the only way a unit's
entry of $D$ can change is for its pre-activation to cross zero. ``The
activation breaks at $z=0$'' means exactly this.

\emph{Frozen pre-activations $z^{(\ell)}_{\sigma,q}$.} Fix $\sigma$ and pretend
every unit below layer $\ell$ has the slope $\sigma$ prescribes, regardless of
the input. Every layer below $\ell$ is then a fixed linear map followed by
multiplication by a fixed diagonal, so the composite is affine and
$z^{(\ell)}_{\sigma,q}(x)=\alpha^\top x+\beta$ for some $\alpha\in\R^d$,
$\beta\in\R$. This is a polynomial of degree at most $1$, and it is
\emph{affine}, not homogeneous: $\beta\neq0$ in general, because the biases
contribute to it. It agrees with the network's realised pre-activation exactly
on the inputs whose pattern below $\ell$ is $\sigma$.

\emph{Zero set, and affine hyperplane.} The \emph{zero set} of
$z^{(\ell)}_{\sigma,q}$ is $\{x\in\R^d: \alpha^\top x+\beta=0\}$ --- the zeros
of that degree-$1$ polynomial, and nothing to do with higher-degree varieties.
When $\alpha\neq0$ this set is an \emph{affine hyperplane}: a translate of a
linear subspace of dimension $d-1$, i.e.\ $\{x:\alpha^\top x=-\beta\}$. It is
closed and Lebesgue-null. When $\alpha=0$ the set is either empty ($\beta\neq0$)
or all of $\R^d$ ($\beta=0$), which is the case the proof handles separately.

\emph{Non-local-constancy set.} A function $g$ is \emph{locally constant at}
$x$ if there is \emph{some} open $U\ni x$ on which $g$ is constant. Its
non-local-constancy set is the set of $x$ where no such $U$ exists. Negating
the quantifiers: $x$ lies in that set iff for \emph{every} open $U\ni x$ there
\emph{exists} $y\in U$ with $g(y)\neq g(x)$ --- the inner quantifier over $y$
is existential, not ``for all''. (Reading it as ``$g(y)\neq g(x)$ for all
$y\in U$'' inverts that quantifier and makes the set look open; it is closed.) We write
$\mathcal B_\ell$ for the non-local-constancy set of the pattern truncated to
layers $1,\dots,\ell$, and $\mathcal B=\mathcal B_L$ for the whole pattern.
\end{clarifybox}

\begin{lemma}[Null switching set]\label{lem:null}
Let $\mathcal B$ be the set of inputs at which the slope diagonal of a network
satisfying (LCS) (Definition~\ref{def:lcs}) is not locally constant --- by
Proposition~\ref{prop:lcs}(ii) the activation breaks only at $z=0$, which is why
the walls below are the zero sets and nothing else. Then $\mathcal B$ is closed and contained in
a finite union of affine hyperplanes; its complement $X_{\mathrm{reg}}$ is
open, dense, and of full Lebesgue measure.
\end{lemma}

\begin{proof}
\emph{$X_{\mathrm{reg}}$ is open, hence $\mathcal B$ is closed.} Let
$x\in X_{\mathrm{reg}}$, so the pattern is constant on some open $U\ni x$. That
same $U$ witnesses local constancy at every one of its points: for $y\in U$ the
set $U$ is an open neighbourhood of $y$ on which the pattern is constant, so
$y\in X_{\mathrm{reg}}$. Hence $U\subseteq X_{\mathrm{reg}}$, and
$X_{\mathrm{reg}}$ is open; $\mathcal B$ is its complement and so is closed.
(The point worth noticing is that one does not construct a new neighbourhood
for each $y$ --- the witness $U$ is reused unchanged.)

\emph{The hyperplane bound.} Induct over layers. For each pattern $\sigma$ and unit
$(\ell,q)$ let $z^{(\ell)}_{\sigma,q}$ be the affine function of $x$ obtained
by freezing the masks below layer $\ell$ to $\sigma$. Let $\mathcal B_{\ell}$
be the non-local-constancy set of the pattern up to layer $\ell$;
$\mathcal B_0=\emptyset$. If $x\notin\mathcal B_{\ell-1}$, the realised
pattern below $\ell$ is constant on a neighbourhood $U$ of $x$, on which the
realised $z^{(\ell)}_q$ equals $z^{(\ell)}_{\sigma,q}$. If
$z^{(\ell)}_{\sigma,q}\equiv0$ then the unit's mask bit is locally constant on
$U$ (an affine function vanishing on an open set vanishes identically, so
there is no partial-vanishing case). Otherwise its zero set is a hyperplane,
off which the bit is locally constant. Hence
$\mathcal B_\ell\subseteq\mathcal B_{\ell-1}\cup\bigcup_{\sigma,q}
\{z^{(\ell)}_{\sigma,q}=0,\ z^{(\ell)}_{\sigma,q}\not\equiv0\}$, a finite
union (patterns that are never realised only enlarge it). For max-pooling, add
the finitely many score-tie hyperplanes.

\emph{Reading off the three conclusions.} Unwinding the induction to
$\ell=L$ puts $\mathcal B=\mathcal B_L$ inside a finite union of sets of the
form $\{z^{(\ell)}_{\sigma,q}=0\}$ with $z^{(\ell)}_{\sigma,q}\not\equiv0$;
each is an affine hyperplane, and the union is finite because there are
finitely many candidate patterns $\sigma$ and finitely many units $q$. That
gives the containment. A hyperplane has Lebesgue measure zero and a finite
union of null sets is null, so $\lambda(\mathcal B)=0$ and
$X_{\mathrm{reg}}=\R^d\setminus\mathcal B$ has full measure. Density follows
from the containment as well: a finite union of hyperplanes has empty interior,
so every ball meets its complement, and $X_{\mathrm{reg}}$ is dense. Openness
was the first paragraph. \qedhere
\end{proof}

\begin{clarifybox}
\textbf{Why this lemma is piecewise-linear and the definition is not.} The
knowledge matrix is defined for \emph{any} activation through the
quotient diagonal \eqref{eq:D-def}, and its row-sum identity needs only that no
pre-activation vanish ($x\in X_{\mathrm{nz}}$) --- nothing about $f(0)$, and
nothing about differentiability. This lemma is a different matter: it is about the set
where the \emph{pattern} is locally constant, and the proof uses that each
$z^{(\ell)}_{\sigma,q}$ is affine, so its zero set is a hyperplane. Both fail
for a strictly nonlinear smooth activation. There $D^{(\ell)}(x)$ takes a
continuum of values and varies continuously with $x$, so the set of inputs with
a locally constant pattern is not of full measure but generically \emph{empty},
and there is no switching set for the lemma to bound. The lemma is therefore
not a technical convenience one could remove by a better definition of
$D^{(\ell)}$: it exists to support the germ identity
(Theorem~\ref{thm:germ}), and it is the germ reading, not the matrix, that is
piecewise-linear in nature. A direct check makes the boundary concrete. At
$x_0=1$ a single $\tanh$ unit with unit weights and the affine map
$x\mapsto f'(1)\,x+\bigl(f(1)-f'(1)\bigr)$ have the \emph{same} germ,
$\bigl(0.419974,\,0.761594\bigr)$, and the same row sum, yet knowledge matrices
$[\,0.761594\mid0\,]$ and $[\,0.419974\mid0.341620\,]$; repeating the
construction with ReLU returns $[\,1\mid0\,]$ for both, as
Theorem~\ref{thm:germ} requires. So for non-PL activations the matrix is not a
function of the germ, and Theorems~\ref{thm:germ}, \ref{thm:maxinv}
and~\ref{thm:complete} are statements about the PL case specifically. Whether
some weaker function-determination survives there is open; we do not claim it.
\end{clarifybox}

\subsection{Proof of Theorem~\ref{thm:germ} (germ identity)}

On a neighbourhood $U$ of $x\in X_{\mathrm{reg}}$ the pattern is constant, so
for every unit $h^{(\ell)}=D^{(\ell)}(x)\,z^{(\ell)}$ on $U$ ($z_q>0$ gives
$\mathrm{ReLU}(z_q)=z_q$; $z_q\le0$ gives $0$). Unrolling,
$\Netxp=\Jwf(x)\,x'+\cwf(x)$ on $U$; differentiate at $x$ and solve for
$c$. For the boundary complement: the same pointwise computation shows
$h^{(\ell)}(x')=D^{(\ell)}(x)\,z^{(\ell)}(x')$ for every $x'$ in the closed
region $\overline{S(x)}$ of $x$'s realised pattern, so the affine map
$A_x(x')=Jx'+c$ agrees with $\Net$ on $\overline{S(x)}$; since $x\in S(x)$
always, $\M(x)\mathbf 1=Jx+c=\Netx$ at every $x$, and for $v$ in the tangent
cone of $S(x)$ at $x$, $\partial^+_v\Net(x)=Jv$ (the one-sided germ from the
inactive side, which is what the convention $\mathbb 1[z>0]$ selects).
\emph{Measured:} probe vs.\ masked product $\le 7\times10^{-15}$; vs.\
finite-difference Jacobian $\le 4\times10^{-11}$; row-sum identity at an
engineered exact-boundary point: $0.0$. \qed

\subsection{Proof of Theorem~\ref{thm:maxinv} (maximal invariance)}

Immediate from Theorem~\ref{thm:germ}: both matrices equal
$[\,D\Net(x)\,\mathrm{diag}(x)\mid \Netx-D\Net(x)x\,]$, a function of the shared germ
alone. The group statement follows because any function-preserving
transformation (any architecture change included) preserves the germ at every
$x\in X_{\mathrm{reg}}$ of both realisations. Sharpness at boundaries: the
identity map and $g(x)=\mathrm{ReLU}(x-1)-\mathrm{ReLU}(1-x)+1\equiv x$
realise the same function, yet at the kink $x=1$ the conventions give
$\M_{\mathrm{id}}(1)=[\,1\mid0\,]$ and $\M_g(1)=[\,0\mid1\,]$ (equal row
sums). \qed

\subsection{Proof of Lemma~\ref{lem:quiver-inv} (quiver-isomorphism invariance) and Proposition~\ref{prop:perm} (permutation scope)}

A quiver isomorphism factors into a per-layer neuron permutation $P_\ell$ and a
positive-diagonal rescaling $T_\ell$ on each hidden layer; we treat the two
generators and compose.

\emph{Permutations.} Let $P_\ell$ permute hidden layer $\ell$, with
$W^{(1)}\!\to P_1W^{(1)}$, $b^{(\ell)}\!\to P_\ell b^{(\ell)}$,
$W^{(\ell)}\!\to P_\ell W^{(\ell)}P_{\ell-1}^{\top}$,
$W^{(L)}\!\to W^{(L)}P_{L-1}^{\top}$, and neuron-wise activations carried with
the units. By induction $z^{(\ell)}(W^{\pi},x)=P_\ell z^{(\ell)}(W,x)$
for every $x$, so the diagonal secant matrices satisfy
$A^{(\ell)}_{\pi}(x)=P_\ell A^{(\ell)}(x)P_\ell^{\top}$ (entrywise ratios
permute; the $0/0\to0$ guard is applied entrywise, hence equivariantly ---
this is where non-equivariant \emph{vector} activations are excluded). The
slope product and the bias accumulation then telescope:
$W^{(L)}P_{L-1}^{\top}\cdot P_{L-1}A^{(L-1)}P_{L-1}^{\top}\cdot
P_{L-1}W^{(L-1)}P_{L-2}^{\top}\cdots=W^{(L)}A^{(L-1)}W^{(L-1)}\cdots$, so
$\KMarg{W^\pi}{x}=\KMwfx$.

\emph{Positive rescalings.} For $T_\ell=\mathrm{diag}(\tau^{(\ell)})$ with
$\tau^{(\ell)}_q>0$, conjugate $W^{(\ell)}\!\to T_\ell W^{(\ell)}T_{\ell-1}^{-1}$,
$b^{(\ell)}\!\to T_\ell b^{(\ell)}$, with the per-neuron activation rescaled to
$\varphi_{\tau}(z)=\tau\,\varphi(z/\tau)$ (for ReLU, $\varphi_\tau=\varphi$).
Then $z^{(\ell)}\!\to T_\ell z^{(\ell)}$ and
$A^{(\ell)}_{\tau}(x)=T_\ell A^{(\ell)}(x)T_\ell^{-1}$, so the inserted factors
$T_\ell^{-1}T_\ell=I$ telescope by the identical argument and
$\KMarg{W^\tau}{x}=\KMwfx$. Composing the two generators gives the lemma.
This recovers the ReLU case of Thms.~4.1--4.2 of the original version.

\emph{Permutation scope (Proposition~\ref{prop:perm}).} The proposition is the
permutation case above, so invariance is immediate; the content is the scope.
The telescoping used that $P_\ell$ commutes with $\varphi$ applied
neuron-wise, which requires elementwise (equivariant) activations; non-equivariant
vector activations break the equivariance of the secant guard above. Moreover,
for activations with $\varphi(0)\neq0$ the row-sum identity itself can fail at
exact $z=0$ points (by $a_q\varphi(0)$; measured example $0.5$), because $P$ no
longer commutes with the constant offset there: a measure-zero caveat the
implementation should carry.
\emph{Measured:} drift $0.0$ at generic and at engineered boundary points,
ReLU and GELU-secant. \qed

\subsection{Proof of Theorem~\ref{thm:complete} (completeness)}

\emph{(i)} $J_{:,i}=\M_{:,i}/x_i$ for $x_i\neq0$ and $c=\M_{:,d+1}$.
\emph{(ii)} $\Net=\M\mathbf 1$ pointwise gives determination on any set; on a
full-measure set, density plus continuity of $\Net$ extends it globally. The germ
at $x$ extends $\Net$ to $\overline{S(x)}$ (proof of Theorem~\ref{thm:germ}),
but no further: $\Psi_1=\mathrm{ReLU}(x-1)$ and $\Psi_3=\mathrm{ReLU}(x)$ have
identical fields ($\equiv[\,0\mid0\,]$) on $(-1,0)$ yet differ at $0.75$ ---
so a field on an open set does not determine $\Net$ even on the region closures
of \emph{other} candidate networks.
\emph{(iii)} For $\tau_q>0$ per hidden unit with weights conjugated,
$z^{(\ell)}\mapsto T_\ell z^{(\ell)}$ and ReLU commutes with positive scalars,
so $a_v\mapsto\tau_v a_v$ while $\Net$ is unchanged.
\emph{(iv)} Write $J=\sum_q D^{(1)}_{qq}\,U_{:,q}\,w_q^{\top}$ with
$U=W^{(L)}D^{(L-1)}\cdots W^{(2)}$. Since $D\Net(x)=J\neq0$, some first-layer
unit $q$ is active with $u:=U_{:,q}\neq0$. Perturb $w_q'=w_q+t e_i$,
$b_q'=b_q-t x_i$: then $z_q'(x)=z_q(x)$ \emph{algebraically}, every hidden
activation at $x$ and $\Netx$ are unchanged, all masks at $x$ are unchanged and
strict (so $x$ stays regular), while the germ moves by the exact rank-one
update $J'-J=u\,(te_i)^{\top}\neq0$ and $c'-c=-t x_i u$, preserving
$J'x+c'=Jx+c$. Two affine maps with different linear parts differ off a
hyperplane, hence on every neighbourhood of $x$; and for any $x_i\neq0$,
column $i$ of $\M$ moves by $t\,x_i\,u\neq0$. (One hidden layer is needed; in
the affine case the same compensation changes nothing.)
\emph{Measured:} activation record equal to $4.4\times10^{-16}$ (2\,ulp),
$\Netx$ equal exactly, rank-one identity $4.7\times10^{-17}$, germ change
$1.9\times10^{-4}\neq0$ at $t=10^{-3}$. \qed

\subsection{Proof of Theorem~\ref{thm:weff} (gradient$\times$input $\oplus$ bias; route equivalence)}

The first $d$ columns of $\M(x)$ are $J_{:,i}x_i=(\partial\Net/\partial x_i)\,x_i$
by the germ identity ($J=D\Net$ at regular $x$, Theorem~\ref{thm:germ}), which is
per-class gradient$\times$input; the bias identity below gives the last column.
The three computation routes (masked product (a), probe (b), autograd (c))
agree, as follows.
(a)$=$germ by Theorem~\ref{thm:germ}. (b): the mask-frozen network is the
affine map $x'\mapsto Jx'+c$; for ReLU the secant $h/z$ equals
$\mathbb 1[z>0]$ ($0/z=0$ for $z<0$; $0/0\to0$ by convention), so probing
$x_ie_i$ through its linear part returns column $J_{:,i}x_i$, and the
zero-input pass through its affine part returns $c$. (c): at regular $x$, $\Net$
is differentiable with $D\Net=J$, and the autograd conventions
($\mathrm{ReLU}'(0)=0$, stored pooling argmax) match the frozen masks even on
the null set. The bias identity: $\partial\Net/\partial b^{(\ell)}
=W^{(L)}D^{(L-1)}\cdots D^{(\ell)}$, and unrolling the bias accumulator gives
$c=\sum_\ell(\partial\Net/\partial b^{(\ell)})\,b^{(\ell)}$.
\emph{Measured} (AlexNet, float64, eval mode): library probe vs.\ autograd
germ $3.3\times10^{-17}$ max-abs (relative $1.3\times10^{-15}$); row sums
$3.5\times10^{-17}$; FullGrad bias identity $6.1\times10^{-16}$. The
network must be in evaluation mode: with dropout active, the saving, probe,
and autograd passes sample different masks and the identity fails at
$\sim10^{-2}$. \qed

\subsection{Proof of Proposition~\ref{prop:contraction} (contraction gap)}

Take the $1$--$2$--$1$ bias-free pair $A$: $w=(1,2)$, $a=(1,1)$ and $B$:
$w=(1,2)$, $a=(1.5,0.75)$. Both realise $\Netx=3\,\mathrm{ReLU}(x)$, so
$\M_A\equiv\M_B$ (measured: $0.0$ on a 401-point grid including the kink). The
per-unit path values $a_qw_q$ are $\{1,2\}$ vs.\ $\{1.5,1.5\}$; any
isomorphism rescales $(w_q,a_q)\mapsto(\tau_qw_q,a_q/\tau_q)$, preserving each
$a_qw_q$, and permutations permute the multiset --- since the multisets
differ, the networks are non-isomorphic and the path-resolved representations
differ at every $x>0$. The identifiability statement on the
\citet{phuong2020functional} class follows by chaining $\Net=\M\mathbf 1$
(field determines the function on a full-measure subset of their domain $Z$,
hence on $Z$ by continuity) with their Theorem~1, quantifiers intact: there
exists $W^{*}$ such that any \emph{general} network of the \emph{same
architecture} agreeing with it on $Z$ is permutation$\times$positive-rescaling
equivalent; the converse direction is Proposition~\ref{prop:perm} plus
rescaling invariance. (Note the $1$--$2$--$1$ example has increasing widths,
consistently outside their architecture class.) \qed

\subsection{Proof of Theorem~\ref{thm:pyth} (visible/invisible decomposition)}

For any $a\in\R^C$ and any $Q'$ with $Q'\mathbf 1=0$:
$\langle a\mathbf 1^{\top},Q'\rangle_F=a^{\top}(Q'\mathbf 1)=0$, so
$\mathcal V=\{a\mathbf 1^{\top}\}$ and $\mathcal N=\{Q:Q\mathbf 1=0\}$ are
orthogonal complements; $A\mapsto(A\mathbf 1)\mathbf 1^{\top}/(d{+}1)$ is the
orthogonal projector onto $\mathcal V$. With
$P=\Delta\Net\,\mathbf 1^{\top}/(d{+}1)$:
$P\mathbf 1=\Delta\Net$, $Q=\Delta\M-P\in\mathcal N$, and
$\|P\|_F^2=\|\Delta\Net\|^2/(d{+}1)$, giving the Pythagorean identity. The
constraint set $\{A:A\mathbf 1=\Delta\Net\}$ is $P+\mathcal N$, so $P$ is its
unique minimum-norm element; equality in the floor iff $Q=0$, i.e.\ constant
rows (per row, Cauchy--Schwarz against $\mathbf 1$). Monotonicity of
$t\mapsto1/\sqrt{(d{+}1)t}$ gives the rank-statistic correspondence (for
medians of even-count samples the interpolated median commutes only up to the
gap between the central order statistics --- below quoting precision in our
tables). Invariances: output rescaling and gauge invariance are immediate from
invariance of $\M$ and linearity; for $x\mapsto Sx$,
$W^{(1)}\mapsto W^{(1)}S^{-1}$ (diagonal $S$): $J\mapsto JS^{-1}$ and
$\mathrm{diag}(Sx)=S\,\mathrm{diag}(x)$, so $J\,\mathrm{diag}(x)$ is
unchanged. Non-invariance under rotations is exhibited numerically (function
preserved to $1.8\times10^{-15}$, $d_M$ changed by up to $57\%$).
\emph{Measured:} identities $\le1.7\times10^{-15}$ over $2{,}400$ random
instances; min-norm never beaten in $10^4$ trials. \qed

\subsection{Proof of Theorem~\ref{thm:within} (within-region anatomy)}

Same strict pattern at $x$ and $y$ puts the segment in one convex region with
common $(J,c)$, so $\M(y)-\M(x)=[\,J\,\mathrm{diag}(y)-J\,\mathrm{diag}(x)
\mid c-c\,]=[\,J\,\mathrm{diag}(\delta)\mid0\,]$ and the column-energy formula
follows. (i): with $a_i=\delta_iJ_{:,i}$,
$\|\sum_ia_i\|^2\le(\sum_i\|a_i\|)^2\le d\sum_i\|a_i\|^2$ (triangle, then
Cauchy--Schwarz), equality iff all $a_i$ equal and nonzero --- attained by
$J=u\mathbf 1_d^{\top}$, $\delta=\mathbf 1_d$. (ii): a single nonzero column
$a=s\,J_{:,i_0}$ gives $d_M=\|a\|=d_\Psi$ and exactly one nonzero column with a
zero bias column. (iii): at most $k$ nonzero columns; the same
Cauchy--Schwarz over the support gives $A\le k$ and the participation bound.
\emph{Measured:} closed form $\le5\times10^{-13}$; cap never violated
($2\times10^4$ trials); one-pixel law exact; $A=1+5.5\times10^{-10}$ in the
finite-step demo. \qed

\subsection{Proof of Theorem~\ref{thm:anatomy} (smooth $+$ crossing)}

\emph{Dyad lemma.} At a transversal single flip of unit $k$ in layer
$\ell^{*}$ at point $z$ on the segment, with $\sigma=\pm1$ the mask change,
the region data jump by $J_{+}-J_{-}=\sigma\,u_kv_k^{\top}$ and
$c_{+}-c_{-}=\sigma\,\gamma_ku_k$, where
$u_k=U^{(\ell^{*}+1)}W^{(\ell^{*}+1)}e_k$,
$v_k^{\top}=e_k^{\top}W^{(\ell^{*})}L^{(\ell^{*})}$, and
$z^{(\ell^{*})}_k(x')=v_k^{\top}x'+\gamma_k$; continuity of $\Net$ across the
wall forces the bias jump, and $v_k^{\top}z+\gamma_k=0$ at the crossing makes
the full dyad's row sums vanish:
$\sigma u_k(v_k^{\top}z+\gamma_k)=0$.
\clarify{Here $U^{(\ell^{*}+1)}$ and $L^{(\ell^{*})}$ denote the downstream and
upstream masked products of the region data shared by the two sides of the
wall, i.e.\ $U^{(\ell^{*}+1)}=W^{(L)}D^{(L-1)}\cdots D^{(\ell^{*}+1)}$ and
$L^{(\ell^{*})}=D^{(\ell^{*}-1)}W^{(\ell^{*}-1)}\cdots D^{(1)}W^{(1)}$ in the
factorisation $J=W^{(L)}D^{(L-1)}\cdots D^{(1)}W^{(1)}$ of
Section~\ref{sec:germ}; so $u_k\in\R^{C}$ is unit $k$'s output-side vector,
$v_k\in\R^{d}$ its input-side vector, and $\gamma_k$ the bias accumulated in
its pre-activation.}

\emph{Telescope.} Write
$\M(y)-\M(x)=\sum_{j=0}^{N}\bigl[\M^{(j)}(z_{j+1})-\M^{(j)}(z_j)\bigr]
+\sum_{j=1}^{N}\bigl[\M^{(j)}(z_j)-\M^{(j-1)}(z_j)\bigr]$, where
$\M^{(j)}(z)=[\,J_j\,\mathrm{diag}(z)\mid c_j\,]$ is the $j$-th region's
matrix evaluated at $z$ (well defined on closures). Each smooth term is
$[\,J_j\,\mathrm{diag}((t_{j+1}-t_j)\delta)\mid0\,]$ by
Theorem~\ref{thm:within}; summing gives $\bar J$. Each jump term is a dyad of
the lemma. Row sums: the smooth part gives $\bar J\delta$, the dyads give $0$,
and $\Nety-\Netx=\int_0^1J(x+t\delta)\,\delta\,dt=\bar J\delta$ by the
fundamental theorem of calculus for the PL map on the segment. The bias
column of the smooth part is zero, so the bias column of $\M(y)-\M(x)$ is
exactly $c(y)-c(x)=\sum_j\sigma_j\gamma_{k_j}u_{k_j}$ --- equivalently, it is
constant on chambers of pair space with fixed crossing combinatorics, which is
the soundness direction of the crossing detector (no crossings $\Rightarrow$
zero bias column). The converse fails: bias-free networks have all
$\gamma_k=0$ (measured: $19$ crossings, bias column exactly $0$), a flipped
unit cut off downstream has $u_k=0$, and multi-crossing cancellations exist.
Genericity (finitely many transversal single crossings) holds off a null set
of pairs by a routine transversality argument over the finite wall complex,
assuming no two units share a wall piece --- exactly duplicated units in
pruned or weight-tied networks violate it and should be screened.
\emph{Hamming.} A unit whose endpoint mask bits differ flips an odd number of
times ($\ge1$); equal bits flip an even number ($\ge0$): $H\le N_{\mathrm{cross}}$
with parity equality, and equality iff no unit flips twice (first-layer walls
are flat, so first-layer units never flip back along a segment; deeper walls
are bent and do --- measured: a hat-shaped network with $N=4$, $H=2$).
\emph{Measured:} full decomposition $\le2.6\times10^{-16}$ over $>1{,}000$
crossings with two independent traversal implementations; bias-column formula
$\le8.2\times10^{-13}$. \qed

\subsection{Proof of Proposition~\ref{prop:dichotomy}}

On the gauge orbit of $W$ (positive rescalings), $\Net$, $d_\Psi$, and $d_M$
are fixed (invariance of $\Net$ and of $\M$), while the penultimate distance
$d_h$ between a fixed pair's features sweeps an unbounded range (rescale one
penultimate unit by $\lambda$: measured $d_h\in[0.44,\,15.3]$ with $\Net$
bitwise fixed). Hence any gauge-invariant function of the stored statistics
$(d_h,d_\Psi)$ is constant in $d_h$ on orbits, i.e.\ uses no penultimate
information. The KM accounting exists because the constraint
$\M\mathbf 1=\Net$ pairs all matrices against the \emph{fixed} vector
$\mathbf 1$; the last-layer relation $\Delta\Net=W^{(L)}\Delta h$ pairs features
against a parameter-dependent map, whose induced ``visible fraction''
$1/(1+\lambda^2)$ in the worked example takes every value in $(0,1)$ across
one orbit. \qed

\subsection{Proof of Theorem~\ref{thm:nogo} and Proposition~\ref{prop:visdrift}}

\emph{No-go.} In $d=C=1$ take
$g(t)=K\,\mathrm{ReLU}(t-x_0)-K\,\mathrm{ReLU}(t-x_0-\varepsilon/K)$ and
$\Net\equiv0$ realised with the same first layer and zeroed second layer. Then
$\sup_t|\Psi-g|=\varepsilon$ exactly; at $x=x_0+\varepsilon/(2K)$ both networks
have masks $(1,0)$; but $J_g(x)=K$, $J_\Psi(x)=0$, so
$\|\M_g(x)-\M_\Psi(x)\|_F=K\sqrt{x^2+x_0^2}\to\infty$ as $K\to\infty$ at fixed
$\varepsilon$ (the construction needs the ramp away from the origin,
$x_0\neq0$). \emph{Measured:} drift $1.4\times10^{2},10^{4},10^{6}$ at
$K=10^{2},10^{4},10^{6}$, $\varepsilon=10^{-2}$ fixed.
\emph{Visible bound.} $\|(\M_\tau(x)-\M(x))\mathbf 1\|=\|\Psi_\tau(W\!,f)(x)-\Netx\|$ is
the row-sum identity; the projection onto constant-row matrices then has norm
$\le\varepsilon/\sqrt{d+1}$. \emph{Conditional bound.} If both networks are
affine on $B(x,r)$ with $\varepsilon_r=\sup_{B(x,r)}\|\Psi_\tau(W\!,f)-\Net\|$, then for
any unit vector $u$, $2r\,\Delta J\,u$ is a difference of differences of
$\Psi_\tau(W\!,f)-\Net$ at $x\pm ru$, so $\|\Delta J\|_{\mathrm{op}}\le\varepsilon_r/r$;
combine $\|\Delta J\,\mathrm{diag}(x)\|_F\le\|x\|_\infty\sqrt{\min(C,d)}\,
\|\Delta J\|_{\mathrm{op}}$ with
$\|\Delta c\|\le\varepsilon_r+\|\Delta J\|_{\mathrm{op}}\|x\|_2$.
\emph{Measured} (exact trust-region computation of $\varepsilon_r$):
$0/120$ violations, minimum slack $0.21$. \qed
